\section{Excess Gaussian energy and intrinsic product complexity}
\label{sec:intrinsic-product-energy}

Define a function-level complexity, allowing all finite architectures,
by
\[
 \mathfrak P(f)=\inf\left\{
   \prod_j\|W^j\|_F:
   f(x)z_w(x)\ge1\ \hbox{on the entire cube}
   \right\}.
\]
The infimum permits arbitrary finite depth and width, retains tanh
hidden layers and a linear readout, and allows biases only through
the augmented input. It is finite, for example by a truth-table
shallow construction. It is also at least one: the network extends
oddly to the full augmented cube $\{-1,1\}^{d+1}$ and has absolute
margin at least one there. This nonconstant odd Boolean extension
has a bichromatic cube edge. Along that edge the product of layer
operator norms, and hence the Frobenius product, must be at least one.

Let $J_{L,N}(f)$ be the infimum of the Gaussian energy $\kappa_N$
over width-$N$, depth-$L$ unit-margin representations of $f$, with
$J_{L,N}(f)=+\infty$ if none exists. This is a deterministic energy
quantity, not a finite-temperature posterior probability.

\begin{theorem}[Intrinsic product complexity characterizes excess energy]
\label{thm:intrinsic-product-energy}
For every $d\ge1$, $L\ge2$ and $N\ge1$,
\[
 J_{L,N}(f)\ge L+\log d+2\log\mathfrak P(f).
\]
Conversely, there are universal polynomials $p,q$ and a universal
constant $c$ such that
\[
 L\ge p(\mathfrak P(f))+c\log(d+1),\qquad
 N\ge q(d,\mathfrak P(f))
\]
imply
\[
 J_{L,N}(f)\le L+c\log(2d\mathfrak P(f)).
\]
The upper-bound witness can be chosen with all physical weights at
most $q(d,\mathfrak P(f))$. The polynomials do not depend on any
supplied representation's depth, width, or coefficient precision.
\end{theorem}

\begin{proof}
The product of the $L$ weighted layer energies is
$d\prod_j\|W^j\|_F^2$. Arithmetic--geometric mean and
$e^u\ge1+u$ therefore give
\[
 \kappa_N\ge L\left(d\prod_j\|W^j\|_F^2\right)^{1/L}
 \ge L+\log d+2\log\mathfrak P(f).
\]
This proves the lower bound without any width or depth restriction.

Choose a finite representation of product norm less than
$2\mathfrak P(f)$. The depth-independent normal form supplies a
unit-margin representation whose depth $D_0$, width $C$, physical
weights and product norm $P_0$ are all polynomial in $\mathfrak P(f)$.
Apply the product-target construction of
Theorem~\ref{thm:product-target-log-padding}. Its width threshold $M$
and transmission parameter $T$ are polynomial in $d,C,D_0,P_0$,
hence polynomial in $d,\mathfrak P(f)$, and do not depend on the
subsequently chosen learner depth. For
$L\ge D_0+\lceil\log T\rceil$, that construction has energy at most
\[
 L+15+32\log T.
\]
If the small-prefix linearization in that proof removes every hidden
activation, its separate halfspace learning conclusion is not needed
here. The collapsed affine map has norm at most $2P_0$ and margin
greater than one. Apply tanh to that map and append readout $2$.
This is a width-one shallow unit-margin representation with weights
polynomial in $P_0$. The shallow scalar-transmission construction
therefore supplies, at every sufficiently large polynomial width and
every $L\ge2+O(\log(d+P_0))$, energy
$L+O(\log(2dP_0))$. It also has polynomially bounded physical weights.
Thus the affine branch has the same claimed upper-bound form.

Its physical weights are bounded by a polynomial width threshold.
Since $\log T=O(\log(d+1)+\log\mathfrak P(f))$, increasing universal
polynomials and constants gives the stated conditions and upper
bound. No attainment of the infimum defining $\mathfrak P(f)$ was
assumed.
\end{proof}

\begin{corollary}[Learning from intrinsic product complexity]
\label{cor:intrinsic-product-learning}
For each fixed $\zeta>0$, the unrestricted Gaussian posterior has a
sample bound polynomial in
\[
 d,\mathfrak P(f),L,\varepsilon^{-1},\log(2/\delta)
\]
at every $L\ge p(\mathfrak P(f))+c\log(d+1)$, simultaneously for
all actual widths $N\ge C_L(f)$ and all
$0<\sigma^2\le N^{-(1+\zeta)}$.
\end{corollary}

\begin{proof}
Use the same depth-independent normal form and apply the all-width
product-target learning theorem. Its intrinsic parameters are all
polynomial in $\mathfrak P(f)$. Its finite-width branch covers the
interval between the actual minimum width and the constructed
polynomial-width witness threshold.
\end{proof}

Thus, once depth and width are sufficiently large, the excess above
the depth baseline is, up to universal constants, the logarithm of
intrinsic product complexity together with the input normalization
cost $\log d$. This does not identify $\mathfrak P(f)$ with ordinary
circuit size. The normal form proves that ordinary Boolean description
size is polynomial in $\mathfrak P(f)$; a converse polynomial bound
for arbitrary small deep circuits remains the central unresolved
representation question. The learning corollary also retains its
displayed learner-depth threshold.
